% Осень26!B21; нумерация внутри столбца без пустых строк.
% Источник: лекция 3, §5 (строка 137)
% Источник: лекция 3, §8 (строка 235)
\definitionticket{18}{Мера Лебега, измеримое по Лебегу множество}
\label{definition:18}

\subsection*{Внешняя мера; измеримость по Каратеодори; схема продолжения}

\begin{definition}

\emph{внешняя мера.} Функция \(\mu:2^X\to\overline{\mathbb{R}}\) (заданная на \textbf{всех} подмножествах \(X\)) называется \textbf{внешней мерой}, если:

\begin{enumerate}
\def\labelenumi{\arabic{enumi}.}
\tightlist
\item
  \(\mu\varnothing=0\);
\item
  \textbf{счётная полуаддитивность}: для любых \(A,A_1,A_2,\dots\subset X\) \[A\subset\bigcup_{i=1}^\infty A_i\quad\Rightarrow\quad \mu A\le\sum_{i=1}^\infty\mu A_i.\]
\end{enumerate}

\textbf{Требует сверки.} В лекции перечислены только эти две аксиомы; в учебнике в определение дополнительно включена неотрицательность \(\mu\ge 0\). «Типичная лингвистическая ситуация: милостивый государь --- не государь; \textbf{внешняя мера --- это не мера}: она не обязана быть счётно аддитивной, и этого, вообще говоря, нет». Отличие от полуаддитивности объёма (теорема 19): там требовалась измеримость всех «персонажей» (принадлежность σ-алгебре/алгебре); здесь все множества доступны.

\end{definition}

\begin{definition}

\emph{измеримость по Каратеодори.} Множество \(A\subset X\) называется \textbf{измеримым относительно внешней меры} \(\mu\), если оно «поддерживает аддитивность» этой меры, т.е. \[\forall B\subset X:\qquad \mu B=\mu(A\cap B)+\mu(B\setminus A).\]

Смысл: множество \(A\) рассекает произвольное «пробное» множество \(B\) на две части --- \(B\cap A\) и \(B\setminus A\); из счётной полуаддитивности неравенство \(\le\) выполнено всегда, а требуется \textbf{равенство}. «Множество \(A\) такое хорошее, что какое \(B\) ни возьми --- даже патологическое, --- оно разбивает его на две части качественно, так что можно написать аддитивность».

\end{definition}

\textbf{Утверждение (без доказательства).} Множества, измеримые относительно внешней меры, образуют \textbf{σ-алгебру}.

Лектор: доказательство существует и устроено «в том же антураже» (объединения, пересечения, предельные переходы), но технически тяжеловесно; «в принципе, можно воспринимать как загадку: надо проверить, что всё пространство измеримо (это, видимо, очевидно), что разность измеримых измерима, и подумать про счётное объединение». \textbf{Требует сверки.} Подробные доказательства в лекции не приводились.

\textbf{Схема лебеговского продолжения (по Каратеодори).} Пусть есть полукольцо \(\mathcal{P}\) из «простых, бедных» множеств и на нём задана \textbf{мера} \(\mu_0\) (пример: ячейки и классический объём --- «объём --- это мера на полукольце ячеек; я это формально не проверил, но сегодня проверю»). Тогда:

\begin{enumerate}
\def\labelenumi{\arabic{enumi}.}
\tightlist
\item
  строим внешнюю меру на всех подмножествах: \[\mu^*(A):=\inf\Bigl\{\sum_{i=1}^\infty\mu_0P_i:\ A\subset\bigcup_{i=1}^\infty P_i,\ P_i\in\mathcal{P}\Bigr\};\] «покрыли множество ячейками, посчитали сумму мер; измельчайте их, чтобы плотнее облегали, --- и перейдите к инфимуму»;
\item
  \textbf{скрытые теоремы} (формулируются, не доказываются):

  \begin{itemize}
  \tightlist
  \item
    \(\mu^*\) действительно внешняя мера;
  \item
    всякое \(P\in\mathcal{P}\) измеримо относительно \(\mu^*\) (т.е. полукольцо содержится в классе измеримых множеств);
  \item
    сужение \(\mu^*\) на полукольцо возвращает исходную меру: \(\mu^*\big|_{\mathcal{P}}=\mu_0\);
  \item
    на σ-алгебре \(\mu^*\)-измеримых множеств внешняя мера \(\mu^*\) является \textbf{мерой}, причём \textbf{полной} («бонус, рекламная продукция нашей фирмы: наша мера --- самая лучшая на свете, она ещё и полная»).
  \end{itemize}
\end{enumerate}

\subsection*{\texorpdfstring{Мера Лебега: определение; точки, счётные множества, \(\mathbb{Q}\)}{Мера Лебега: определение; точки, счётные множества, \textbackslash mathbb\{Q\}}}

\begin{definition}

\emph{мера Лебега.} \textbf{Мера Лебега} \(\lambda=\lambda_m\) в \(\mathbb{R}^m\) --- результат применения теоремы о продолжении меры (§6, «продолжение по Каратеодори») к классическому объёму \(\mu\) на полукольце ячеек \(\mathcal{P}^m\): \[(\mathcal{P}^m,\ \mu)\quad\longmapsto\quad(\mathcal{M}^m,\ \lambda),\qquad \lambda\big|_{\mathcal{P}^m}=\mu .\]

\end{definition}

\textbf{Лебеговская σ-алгебра} \(\mathcal{M}^m\) --- σ-алгебра измеримых множеств из теоремы о продолжении (область определения \(\lambda\); при необходимости указывать размерность --- пишут \(\lambda_m\)). Множества \(A\in\mathcal{M}^m\) называются \textbf{измеримыми по Лебегу}.

«В результате продолжения мы приобретаем чёрный ящик: σ-алгебру измеримых множеств. Из чего она состоит, я пока не знаю --- это главный секрет; косвенное описание: то, что хорошо разбивает внешнюю меру. Будем присматриваться и добывать свойства».

\textbf{Свойство 1 (мелкие множества).}

\begin{itemize}
\tightlist
\item
  \textbf{Точка измерима и имеет меру ноль}: \(\lambda\{x\}=0\) для \(x\in\mathbb{R}^m\). \emph{Измеримость}: точка --- пересечение уменьшающихся ячеек («берёте ячейку, ячейку поменьше, ещё меньше --- пересечение будет вот этот уголочек»), т.е. счётное пересечение измеримых множеств, а σ-алгебра замкнута относительно счётных пересечений. \emph{Мера ноль}: по монотонности \(\lambda\{x\}\) не больше меры любой такой ячейки; кубик со стороной \(r\) имеет меру \(r^m\), значит \(\lambda\{x\}\le r^m\) для любого \(r>0\), т.е. \(\lambda\{x\}=0\). («Кто бы мог выдумать!»)
\item
  \textbf{Счётное множество имеет меру ноль}: оно --- счётное дизъюнктное объединение точек, и по счётной аддитивности \(\lambda=\sum 0=0\).
\item
  В частности \(\lambda_1(\mathbb{Q})=0\). «Между прочим, ничего так: \(\mathbb{Q}\) геометрически плотное на всей оси, и тем не менее мера ноль».
\end{itemize}
